(1) 利用 \([a^\dagger a,a]=-a\)、\([\sigma_+\sigma_-,\sigma_+]=\sigma_+\) 及其厄米共轭，可得 \([N,a\sigma_+]=[N,a^\dagger\sigma_-]=0\)，从而 \([H,N]=0\)。固定 \(N=n+1\) 时，
\[
H_n=\hbar\omega_c(n+\tfrac12)\mathbb I+
\frac{\hbar}{2}\begin{pmatrix}\Delta&2g\sqrt{n+1}\\2g\sqrt{n+1}&-\Delta\end{pmatrix}.
\]
记 \(G_n=g\sqrt{n+1}\)、\(\Omega_n=\sqrt{\Delta^2+4G_n^2}\)，定义 \(0<\theta_n<\pi/2\) 满足
\(\cos2\theta_n=\Delta/\Omega_n\)、\(\sin2\theta_n=2G_n/\Omega_n\)。解二阶特征多项式得到
\[
\boxed{E_{n,\pm}=\hbar\omega_c(n+\tfrac12)\pm\tfrac{\hbar\Omega_n}{2}},\quad
|n,+\rangle=\cos\theta_n|e,n\rangle+\sin\theta_n|g,n+1\rangle,\quad
|n,-\rangle=-\sin\theta_n|e,n\rangle+\cos\theta_n|g,n+1\rangle.
\]
唯一 \(N=0\) 态 \(|g,0\rangle\) 与交换项完全解耦，其精确能量为 \(-\hbar\omega_a/2\)；不能把它误当作 \(n=-1\) 的双态块。

(2) 去掉公共相位后，块内无迹哈密顿量为 \(h_n/\hbar=(\Delta\sigma_z+2G_n\sigma_x)/2\)，且 \((h_n/\hbar)^2=\Omega_n^2\mathbb I/4\)。故指数矩阵可直接求得
\[
|\psi(t)\rangle=e^{-i\omega_c(n+1/2)t}
\left[\left(\cos\frac{\Omega_nt}{2}-i\frac{\Delta}{\Omega_n}\sin\frac{\Omega_nt}{2}\right)|e,n\rangle
-i\frac{2G_n}{\Omega_n}\sin\frac{\Omega_nt}{2}|g,n+1\rangle\right].
\]
于是
\[
\boxed{P_{g,n+1}(t)=\frac{4g^2(n+1)}{\Delta^2+4g^2(n+1)}
\sin^2\!\frac{\Omega_nt}{2}},\qquad
P_{\max}=\frac{4G_n^2}{\Omega_n^2}.
\]
交换概率使 \(\langle a^\dagger a\rangle=n+P\)、\(\langle\sigma_+\sigma_-\rangle=1-P\)，故 \(\langle N\rangle=n+1\) 不变。更细致地，裸哈密顿量 \(H_0=\hbar\omega_ca^\dagger a+\hbar\omega_a\sigma_z/2\) 的期望改变量为 \(-\hbar\Delta P\)；交换项 \(V=\hbar g(a\sigma_++a^\dagger\sigma_-)\) 的期望为 \(+\hbar\Delta P\)，因此完整能量严格守恒，失谐时不能只检查裸能量。共振时 \(P=\sin^2(g\sqrt{n+1}\,t)\) 可完全交换；当 \(|\Delta|\gg2G_n\) 时 \(P_{\max}\simeq4G_n^2/\Delta^2\ll1\)。

(3) 相干态展开为 \(|\alpha\rangle=e^{-\bar n/2}\sum_{n=0}^\infty\alpha^n|n\rangle/\sqrt{n!}\)。由于不同激发数块正交，原子反转 \(W=P_e-P_g\) 是
\[
\boxed{W(t)=e^{-\bar n}\sum_{n=0}^\infty\frac{\bar n^n}{n!}
\cos\!\left(2g\sqrt{n+1}\,t\right)}.
\]
令 \(n=\bar n+\delta n\)，在泊松分布的主要权重 \(|\delta n|=O(\sqrt{\bar n})\) 内展开相位：
\(2g\sqrt{n+1}\,t=2g\sqrt{\bar n+1}\,t+
gt\,\delta n/\sqrt{\bar n+1}+O(gt/\sqrt{\bar n})\)（早期 \(gt=O(1)\)）。
利用泊松方差 \(\operatorname{Var}n=\bar n\) 及中心极限定理，得到
\[
\boxed{W(t)\simeq
\cos\!\left(2g\sqrt{\bar n+1}\,t\right)
\exp\!\left[-\frac{g^2t^2\bar n}{2(\bar n+1)}\right]},\qquad
t_{\rm col}\simeq\frac{\sqrt2}{g}\quad(\bar n\gg1),
\]
其中塌缩时间定义为包络降至 \(e^{-1}\)。相邻 \(n\) 的相位差近似
\(2gt(\sqrt{n+2}-\sqrt{n+1})\simeq gt/\sqrt{\bar n+1}\)，首次近似重合为 \(2\pi\)，故
\[
\boxed{t_{\rm rev}\simeq\frac{2\pi\sqrt{\bar n+1}}g}.
\]
平方根频率并不等间隔，二阶展开导致复苏峰展宽和不完全重相位，不能宣称严格周期。

(4) 令 \(\chi=g^2/\Delta\)（带符号）。在 \(|\Delta|\gg2G_n\) 时，
\(\Omega_n=|\Delta|+2G_n^2/|\Delta|+O(G_n^4/|\Delta|^3)\)，结合对应缀饰态与裸态的连续对应，\(|e,n\rangle\) 的二阶能移为 \(+\hbar G_n^2/\Delta\)，\(|g,n+1\rangle\) 为 \(-\hbar G_n^2/\Delta\)。因此
\[
\boxed{H_{\rm eff}=\hbar\omega_ca^\dagger a+
\frac{\hbar\omega_a}{2}\sigma_z+
\hbar\chi(a^\dagger a+\tfrac12)\sigma_z+
\frac{\hbar\chi}{2}\mathbb I}
\]
在占据显著的光子数范围内准确至 \(O(\hbar g^4(n+1)^2/|\Delta|^3)\)；最后一项保证 \(|g,0\rangle\) 的能量仍为 \(-\hbar\omega_a/2\)。在该有效哈密顿量下，初始原子叠加态的非对角元为
\[
\boxed{\rho_{eg}(t)=\tfrac12e^{-i(\omega_a+\chi)t}
\exp\!\left[\bar n(e^{-2i\chi t}-1)\right]},\qquad
\boxed{2|\rho_{eg}(t)|=\exp[-2\bar n\sin^2(\chi t)]}.
\]
第二个指数来自两种原子态对应的相干腔场相位重叠，而非不可逆耗散。短时 \(|\chi t|\ll1\) 时包络为 \(\exp(-2\bar n\chi^2t^2)\)，按 \(e^{-1}\) 定义的色散退相干时间为 \(t_\varphi=1/(|\chi|\sqrt{2\bar n})\)（\(\bar n>0\)）；有效模型在 \(t=\pi/|\chi|\) 完全复相干。实际精确模型的高阶色散项、有限寿命等会限制远时复相干，且 \(\bar n=0\) 时本机制没有退相干。
